% =============================================================================
% INF-221 Algoritmos y Complejidad -- Tarea 1 -- Semestre 2026-2
% Autor: Kurt Wiederhold   Rol: 202473528-4
% Seccion: Entorno experimental (extension maxima: 1/2 pagina)
%
% Los dos archivos que siguen definen, como macros \newcommand, TODOS los
% numeros que la prosa del informe cita (razones, factores, tiempos puntuales).
% Los generan los scripts plot_generator.py a partir de los CSV de mediciones,
% igual que las tablas y las figuras, de modo que el texto no contiene ningun
% valor escrito a mano y no queda desincronizado si se repite el experimento.
% Se cargan aqui porque esta es la primera seccion que los usa.
% =============================================================================
% =============================================================================
% INF-221 Algoritmos y Complejidad -- Tarea 1 -- Semestre 2026-2
% Autor: Kurt Wiederhold   Rol: 202473528-4
%
% MACROS NUMERICAS. ARCHIVO GENERADO AUTOMATICAMENTE por code/sorting/scripts/plot_generator.py
% a partir de data/measurements/sorting_measurements.csv. No editar a mano:
% se regenera con `make plots`.
% =============================================================================
\newcommand{\SortDomainFactorPatience}{9{,}4}
\newcommand{\SortDomainFactorStd}{3{,}1}
\newcommand{\SortExpMax}{1{,}22}
\newcommand{\SortExpMin}{0{,}96}
\newcommand{\SortExpNlogn}{1{,}14}
\newcommand{\SortNMaxMiB}{38{,}1}
\newcommand{\SortNumIncorrect}{0}
\newcommand{\SortNumRows}{2\,880}
\newcommand{\SortOrderGainMerge}{3{,}4}
\newcommand{\SortOrderGainQuick}{4{,}0}
\newcommand{\SortOrderGainStd}{4{,}6}
\newcommand{\SortOthersWorstBestMax}{4{,}6}
\newcommand{\SortOthersWorstBestMin}{3{,}4}
\newcommand{\SortPatAscMs}{1\,063}
\newcommand{\SortPatDescMs}{82{,}5}
\newcommand{\SortPatMemWorst}{156,5~MiB}
\newcommand{\SortPatMemWorstOverArray}{4{,}1}
\newcommand{\SortPatPilesRandom}{6\,325}
\newcommand{\SortPatPilesTheory}{6{,}3}
\newcommand{\SortPatRandMs}{2\,469}
\newcommand{\SortPatRandOverAsc}{2{,}3}
\newcommand{\SortPatWorstBest}{29{,}9}
\newcommand{\SortRatioMergeStd}{1{,}55}
\newcommand{\SortRatioPatienceStd}{4{,}61}
\newcommand{\SortRatioQuickStd}{1{,}22}
\newcommand{\SortStdDescMs}{115{,}2}
\newcommand{\SortStdNsFactor}{4{,}0}
\newcommand{\SortStdNsMid}{2{,}03}
\newcommand{\SortStdNsSmall}{0{,}51}
\newcommand{\SortStdTenNs}{34}
\newcommand{\SortWorstOverBest}{4{,}6}

% =============================================================================
% INF-221 Algoritmos y Complejidad -- Tarea 1 -- Semestre 2026-2
% Autor: Kurt Wiederhold   Rol: 202473528-4
%
% MACROS NUMERICAS. ARCHIVO GENERADO AUTOMATICAMENTE por
% code/matrix_multiplication/scripts/plot_generator.py a partir de
% data/measurements/matrix_multiplication_measurements.csv. No editar a mano:
% se regenera con `make plots`.
% =============================================================================
\newcommand{\MatCtrlDensaMs}{3\,279}
\newcommand{\MatCtrlDiagonalMs}{3\,359}
\newcommand{\MatCtrlDispersaMs}{3\,168}
\newcommand{\MatCtrlShiftPct}{25}
\newcommand{\MatCutoff}{64}
\newcommand{\MatExpNaive}{3{,}525}
\newcommand{\MatExpStrassen}{3{,}066}
\newcommand{\MatExpStrassenPts}{2}
\newcommand{\MatMemRatio}{4{,}0}
\newcommand{\MatNMaxLog}{10}
\newcommand{\MatNMidLog}{8}
\newcommand{\MatNaiveMaxMs}{2\,630}
\newcommand{\MatNaiveMemMax}{8,0~MiB}
\newcommand{\MatNaiveMidMs}{16{,}87}
\newcommand{\MatNumIncorrect}{0}
\newcommand{\MatNumRows}{1\,440}
\newcommand{\MatRsqNaiveMin}{0{,}998}
\newcommand{\MatSpeedupMax}{5{,}67}
\newcommand{\MatSpeedupMid}{2{,}55}
\newcommand{\MatStrassenMaxMs}{464}
\newcommand{\MatStrassenMemMax}{31,9~MiB}
\newcommand{\MatStrideEquivMs}{521}
\newcommand{\MatStrideEquivSpeedup}{1{,}12}
\newcommand{\MatStrideFactor}{6{,}1}
\newcommand{\MatStrideN}{1\,024}
\newcommand{\MatStrideNeighHi}{1\,040}
\newcommand{\MatStrideNeighLo}{1\,000}
\newcommand{\MatStrideNonPowMax}{0{,}50}
\newcommand{\MatStrideNonPowMin}{0{,}31}
\newcommand{\MatStrideNonPowRatio}{1{,}6}
\newcommand{\MatStrideNsNeighMax}{0{,}50}
\newcommand{\MatStrideNsNeighMin}{0{,}47}
\newcommand{\MatStrideNsPow}{2{,}96}
\newcommand{\MatStridePowRatio}{3{,}0}
\newcommand{\MatStrideRows}{54}
\newcommand{\MatSweepBestCut}{16}
\newcommand{\MatSweepBestMs}{4{,}28}
\newcommand{\MatSweepBestOverPure}{23{,}5}
\newcommand{\MatSweepDefaultMs}{6{,}64}
\newcommand{\MatSweepN}{256}
\newcommand{\MatSweepNLog}{8}
\newcommand{\MatSweepNaiveOverBest}{3{,}94}
\newcommand{\MatSweepPureMs}{100{,}6}
\newcommand{\MatSweepPureOverNaive}{6{,}0}
\newcommand{\MatSweepRows}{11\,988}
\newcommand{\MatTypeCheckRows}{18}
\newcommand{\MatTypeDispCtrl}{6}
\newcommand{\MatTypeDispMain}{6}


\textbf{Plataforma.} AMD Ryzen 7 7700 (8 núcleos Zen~4, 3{,}8~GHz), 32~GiB
DDR5, cachés de 32~KiB (L1d, 8 vías), 1~MiB (L2) por núcleo y 32~MiB (L3);
Windows~11 Pro 24H2; \texttt{g++}~16.1.0 (MinGW-w64), opciones
\texttt{-std=c++17} y \texttt{-O2}, sin \texttt{-march=native} para no atar el
binario a este procesador~\cite{hales2017replicability}; Python~3.12.10,
NumPy~2.5.2 y Matplotlib~3.11.1.

\textbf{Datos.} Exactamente los del anexo~A: $72$ arreglos y $72$ pares de
matrices ($4$ tamaños $\times$ $3$ tipos $\times$ $2$ dominios $\times$ $3$
muestras), producidos sin modificación por los generadores entregados
(\texttt{array\_generator.py}, \texttt{matrix\_generator.py}).

\textbf{Metodología.} Se cronometra sólo la llamada al algoritmo, con
\texttt{steady\_clock}; la E/S y la verificación quedan fuera porque para
$n=10^{7}$ leer el archivo tarda más que ordenarlo. Cada caso se repite 15, 7 o
3 veces según su tamaño; se reporta la \emph{mediana}, robusta frente a los
valores atípicos altos del planificador~\cite{mcgeoch2012}. Cuando una
ejecución dura menos que la resolución del reloj ($\approx$ 100~ns) se
cronometra un lote preparado de antemano y se divide el total. La memoria
auxiliar se mide reemplazando los operadores globales
\texttt{new}/\texttt{delete} por versiones instrumentadas. Toda salida se
verifica contra una referencia independiente: las \SortNumRows{} mediciones
de ordenamiento, las \MatNumRows{} de matrices, las \MatSweepRows{} del
barrido de umbral y las \MatTypeCheckRows{}+\MatStrideRows{} de los dos
controles resultaron correctas. Figuras, tablas y toda cifra citada en el
texto se generan desde los CSV con los scripts \texttt{plot\_generator.py}.
